% Editable transcription; verify reused text against the original PDF.
\section{A short history of modern AI}

One constant in the short but rich history of AI is the tension between
two main goals.
The first goal is to understand (human) intelligence and the second is
to let computers do information processing tasks.
Chess is an exemplary example.
The earliest attempts to build computer chess programs were inspired
by studies on human chess thinking (de Groot, 1946).
However, the grand steps were only made when approaches were developed
that made use of the unique strengths of computers.
Deep Blue, the program that beat Gary Kasparov, the world champion in
1996, was a typical cognitive wheel.
It applied advanced tree search algorithms that allowed deep search
of positions based on a rather simple evaluation of resulting positions.
Deep Blue mainly told us how humans do not play chess, as it depended
on a brute force computing approach, evaluating 200 million positions
per second, that humans are particularly bad at.
Humans are much better at evaluating positions (Gobet \& Chassy, 2009).
How chess players think is still largely a mystery
(Gobet \& Charness, 2018), but chess involves both high-level conceptual
processing and low-level recognition of familiar patterns
(Lane \& Chang, 2018; Van Harreveld, Wagenmakers, \& Van Der Maas, 2007;
see Blanch (2020) for an overview of research on individual differences
in chess).

The renewed interest in AI among psychologists is based on novel
techniques that have also revolutionized computer chess in the last
5 years.
AlphaZero chess (Silver et al. 2018) combines a deep learning network
that values positions and predicts next moves, with a reinforcement
learning system (explained in more detail below).
The amazing thing is that AlphaZero initially performs badly, but the
underlying systems help each other to learn and improve during millions
of self-played games.
This approach radically departs from earlier, non-learning, programs
where brute force search and built-in indexes of openings and endgames
were the key to success.
The open-source variant of AlphaZero chess, Leela Chess Zero, currently
dominates chess rating lists and is praised for her human-like play
(including occasional incomprehensible blunders).
So, it could be said that the initial divergence between the two goals
of AI (cognition vs application) has diminished given the latest
spectacular developments in computer chess.

Chess is not the only impressive example of recent successes in AI.
In the last 10 years we have seen breakthroughs in areas such as
natural language (NLP) processing, computer vision, robotics, and more.
In the field of NLP, machine translation, AI assistants
(like Alexa, Google and Siri), and automated closed captions in videos
have become commonplace in our lives.
These have evolved from AI systems that obtained historical milestones
such as IBM's Watson question answering system defeating human
champions at the popular trivia TV show Jeopardy in 2011 and Microsoft's
speech recognition software successfully transcribing conversational
telephone speech on the widely benchmarked Switchboard task in 2017.
More recently, instead of transcribing or interpreting human text or
speech, NLP models have been developed to generate original text and
speech.
A particularly impressive example is the recent GPT-3 model, which was
able to write a complete and coherent article in the Guardian
(GPT-3, 2020) arguing that robots will not destroy humanity, a fear held
by several well-known people in science and industry.

The advances in computer vision are similarly impressive.
Perhaps the most well-known examples are in the domain of object and
face recognition, where Google Lens (available on most Android
smartphones) is an excellent example of object recognition and
Facebook's automatic face tagging system uses the power of face
recognition.
AI powered vision systems have also been applied in other, more
practical domains, such as medicine, in which they are used to assist
medical doctors in disease diagnosis and pharmacologists with
treatment development.
For example, AI models are used to detect cancer from medical scans
and to assist the development of novel medical drugs
(Kaul, Enslin, \& Gross, 2020).

Another field that has seen remarkable developments is robotics.
While robots before the turn of the century were mostly limited to
simple repetitive tasks, modern AI powered robots are able to perform
surgery, have been used in nursing homes to assist with medical
documentation and combat loneliness among their residents, and
commercially available fully self-driving cars are expected to arrive
at the market within 1 or 2 years (Siciliano \& Khatib, 2016).
The developments are spectacular and worrisome at the same time.
The study of the societal consequences of AI is becoming a research
area in itself (Dubhashi \& Lappin, 2017;
Jobin, Ienca, \& Vayena, 2019; Rudin, 2019).
